\documentclass{ctexart}
\usepackage{amsmath,amssymb,amsthm}
\DeclareMathOperator{\sgn}{sgn}
\usepackage{geometry}
\geometry{margin=1in}

\newtheorem{theorem}{定理}[section]
\newtheorem{proposition}[theorem]{命题}
\newtheorem{definition}[theorem]{定义}
\newtheorem{remark}[theorem]{注记}
\newtheorem{assumption}[theorem]{假设}
\newtheorem{lemma}[theorem]{引理}

\title{单目标球形围栏控制 (Spherical Fencing Control of a Single Target)}

\begin{document}

\maketitle

\section{问题描述 (Problem Description)}

考虑 $n$ 个同质智能体 (homogeneous agents) 在 $N$ 维欧几里得空间 (Euclidean space) 中运动，$N \geq 2$。
我们设计协同控制律 (cooperative control laws)，使得智能体集体构成一个半径为 $r > 0$ 的球形编队 (spherical formation)，将目标围在其中。

令 $p_i \in \mathbb{R}^N$ 表示智能体 $i$ 的位置，$p_0 \in \mathbb{R}^N$ 表示目标的位置。
\begin{definition}[球形围栏控制 (Spherical fencing control)]
称目标被智能体\textit{球形围控 (spherically fenced)}，如果有
\begin{equation}
\lim_{t \to \infty} \|p_i(t) - p_0(t)\| = r ,
\quad 
\lim_{t \to \infty} \|\bar{p}(t) - p_0(t)\| = 0
\end{equation}
其中 $\bar{p}=1/n \sum_{i=1}^{n} p_i$ 是 $n$ 个智能体的中心。
\end{definition}

定义 $\sgn(x) = x/\|x\|$（当 $\|x\| \neq 0$ 时），以及 $\sgn(x) = \mathbf{0}$（当 $\|x\| = 0$ 时）。
全文假设每个智能体初始时与目标处于不同位置，即对所有 $i$ 有 $p_i(0) \neq p_0(0)$，且初始相对位置集 $\{p_i(0) - p_0(0)\}_{i=1}^{n}$ 张成 (span) $\mathbb{R}^N$。

\section{智能体动力学 (Agent Dynamics)}

智能体由一阶积分器 (first-order integrators) 控制：
\begin{equation}
\dot{p}_i = u_i \in \mathbb{R}^N, \quad i = 1,2,\dots,n
\label{eq:agent_dyn}
\end{equation}

\section{控制律设计 (Control Law Design)}

定义每个智能体到半径为 $r$ 的球面上的径向投影 (radial projection)：
\begin{equation}
    x_i = r\,\sgn(p_i - p_0) \in \mathbb{R}^N,
\end{equation}
以及智能体 $i$ 与 $j$ 之间的切向耦合向量 (tangential coupling vector)：
\begin{equation}
    d_{ij} = x_j - \frac{\langle x_i, x_j \rangle}{r^2}\, x_i \in \mathbb{R}^N.
\end{equation}
由构造可知，$d_{ij}$ 与 $x_i$ 正交，即 $\langle d_{ij}, x_i \rangle = 0$，因此它位于球面在 $x_i$ 处的切空间 (tangent space) 中。
耦合向量满足：当 $x_i = x_j$ 或 $x_i = -x_j$ 时，$d_{ij} = \mathbf{0}$；其模长为
\begin{equation}
\|d_{ij}\| = \frac{\|x_i - x_j\| \sqrt{4r^2 - \|x_i - x_j\|^2}}{2r}.
\end{equation}

\begin{definition}[智能体间作用力 (Inter-agent Force)]
智能体 $j$ 施加于智能体 $i$ 的交互作用力 (interaction force) 定义为
\begin{equation}
f_{ij} = \left(k_a - \frac{k_r}{\|x_i - x_j\|^2}\right) d_{ij} \in \mathbb{R}^N,
\end{equation}
其中 $k_a > 0$ 为吸引增益 (attraction gain)，$k_r > 0$ 为排斥增益 (repulsion gain)，$d_{ij}$ 为切向耦合向量。
\end{definition}

\begin{remark}[奇异构型 (Singular Configurations)]
当 $x_i = x_j$（两智能体位于同一投影点）时，有 $d_{ij} = \mathbf{0}$，作用力消失，故排斥机制 (repulsive mechanism) 不起作用。
这对应于退化平衡点 (degenerate equilibrium)；实际上，任何微小扰动都会打破对称性并恢复作用力。
类似地，当 $x_i = -x_j$（对径点，diametrically opposite）时，$d_{ij} = \mathbf{0}$，此时也没有切向耦合作用。
\end{remark}

\subsection*{通信图 (Communication Graph)}

令 $\mathcal{G} = (\mathcal{V}, \mathcal{E})$ 为无向图 (undirected graph)，顶点集 (vertex set) 为 $\mathcal{V} = \{1,\dots,n\}$，边集 (edge set) 为 $\mathcal{E}$。
智能体 $i$ 的邻居集为 $\mathcal{N}_i = \{j : (i,j) \in \mathcal{E}\}$，$w_{ij} > 0$ 为边权重 (edge weight)（$j \in \mathcal{N}_i$），且 $w_{ii} = 0$。

\begin{assumption}[图连通性 (Graph Connectivity)]
通信图 $\mathcal{G}$ 是连通的 (connected)。
\end{assumption}

\subsection*{控制律 (Control Law)}

智能体 $i$ 的整体控制输入为
\begin{equation}
u_i =
-k_o\bigl(\|p_i - p_0\| - r\bigr)\,\sgn(p_i - p_0)
+ \sum_{j \in \mathcal{N}_i} w_{ij} \left(k_a - \frac{k_r}{\|x_i - x_j\|^2}\right) d_{ij}
+ \dot{p}_0,
\label{eq u}
\end{equation}
其中 $k_o > 0$ 为径向跟踪增益 (radial tracking gain)。注意第一项可等价地写为 $-\frac{k_o}{r}(\|p_i - p_0\| - r)\,x_i$，即沿径向方向的带符号距离误差之比例项。
控制参数 $k_o, k_r, k_a > 0$ 需满足
\begin{equation}
\frac{k_r}{k_a} \geq \frac{2n\,r^2}{n-1}.
\label{eq:gain_condition}
\end{equation}

\begin{remark}[控制律的结构 (Structure of the Control Law)]
控制输入 \eqref{eq u} 由三项组成：
\begin{enumerate}
    \item \emph{径向跟踪 (Radial tracking)}：$-k_o(\|p_i - p_0\| - r)\,\sgn(p_i - p_0)$ 通过沿径向的带符号距离误差，驱动每个智能体趋向半径为 $r$ 的球面。
    \item \emph{智能体间交互 (Inter-agent interaction)}：邻居节点 $f_{ij}$ 的加权和，通过切向力 (tangential forces) 使智能体在球面上均匀分布。
    \item \emph{前馈 (Feedforward)}：$\dot{p}_0$ 补偿目标的运动。
\end{enumerate}
\end{remark}

\section{稳定性分析 (Stability Analysis)}

\subsection*{误差动力学 (Error Dynamics)}

定义相对位置误差：
\begin{equation}
e_i = p_i - p_0, \quad i = 1,2,\dots,n
\end{equation}
以及径向误差 (radial error)：
\begin{equation}
r_i = \|e_i\|, \quad i = 1,2,\dots,n.
\end{equation}

$e_i$ 的动力学为：
\begin{equation}
\dot{e}_i = u_i - \dot{p}_0 = -k_o(r_i - r)\,\sgn(e_i) + \sum_{j \in \mathcal{N}_i} w_{ij} f_{ij}.
\label{eq: error_dynamics}
\end{equation}

径向误差动力学可推导为：
\begin{equation}
\dot{r}_i = \frac{e_i^T \dot{e}_i}{\|e_i\|} = -k_o(r_i - r) + \frac{e_i^T}{r_i} \sum_{j \in \mathcal{N}_i} w_{ij} f_{ij}.
\label{eq: radial_dynamics}
\end{equation}

\subsection*{辅助结果 (Auxiliary Results)}

\begin{lemma}[作用力的性质 (Properties of the Force)]
对球面 $\mathrm{S}^{(N-1)}$（半径为 $r$）上任意两点 $x_i, x_j$，作用力 $f_{ij}$ 满足：
\begin{enumerate}
    \item $\langle f_{ij}, x_i \rangle = 0$（与球面相切，tangential to the sphere），
    \item 当 $x_i = x_j$ 或 $x_i = -x_j$ 时，$f_{ij} = \mathbf{0}$，
    \item 当 $\|x_i - x_j\| \to 0$ 时，$\|f_{ij}\| \to \infty$。
\end{enumerate}
\end{lemma}

\begin{proof}
由 $f_{ij}$ 的定义及 $\langle d_{ij}, x_i \rangle = 0$，性质 (1) 直接成立。
性质 (2) 和 (3) 由 $d_{ij}$ 的定义和排斥项 (repulsive term) 可得。
\end{proof}

\begin{lemma}[球面约束 (Sphere Constraint)]
在控制律 \eqref{eq u} 下，若对所有 $i$ 有 $r_i(0) = r$，则对所有 $t \geq 0$ 有 $r_i(t) = r$。
\end{lemma}

\begin{proof}
定义 $V_i = \frac{1}{2}(r_i^2 - r^2)$。则：
\begin{equation}
\dot{V}_i = r_i \dot{r}_i = -k_o(r_i - r)r_i + e_i^T \sum_{j \in \mathcal{N}_i} w_{ij} f_{ij}.
\end{equation}
由于 $\langle f_{ij}, x_i \rangle = 0$ 且 $e_i = r_i \sgn(e_i)$，有 $e_i^T f_{ij} = 0$。
因此，$\dot{V}_i = -k_o(r_i - r)r_i$。
若 $r_i(0) = r$，则 $V_i(0) = 0$，且当 $r_i = r$ 时 $\dot{V}_i = 0$，故对所有 $t$ 有 $r_i(t) = r$。
\end{proof}

\subsection*{主要稳定性证明 (Main Stability Proof)}

\begin{theorem}[稳定性 (Stability)]
在假设 1 下，考虑智能体动力学 \eqref{eq:agent_dyn}、控制输入 \eqref{eq u} 以及增益条件 \eqref{eq:gain_condition}。
球形围栏控制问题得以解决：智能体收敛到以目标为中心的半径为 $r$ 的球形编队，即
\begin{equation*}
\lim_{t \to \infty} \|p_i(t) - p_0(t)\| = r, \quad
\lim_{t \to \infty} \|\bar{p}(t) - p_0(t)\| = 0.
\end{equation*}
\end{theorem}

\begin{proof}
我们分几步进行。

\textbf{步骤 1：径向收敛 (Radial convergence)。}
定义候选李雅普诺夫函数 (Lyapunov function candidate)：
\begin{equation}
V_1 = \sum_{i=1}^{n} \frac{1}{2}(r_i - r)^2.
\end{equation}
它沿 \eqref{eq: radial_dynamics} 轨线的导数为：
\begin{equation}
\dot{V}_1 = \sum_{i=1}^{n} (r_i - r)\dot{r}_i = -k_o \sum_{i=1}^{n} (r_i - r)^2 + \sum_{i=1}^{n} (r_i - r) \frac{e_i^T}{r_i} \sum_{j \in \mathcal{N}_i} w_{ij} f_{ij}.
\end{equation}
由于 $\langle f_{ij}, x_i \rangle = 0$ 且 $e_i = r_i \sgn(e_i)$，有 $e_i^T f_{ij} = 0$。
因此，
\begin{equation}
\dot{V}_1 = -k_o \sum_{i=1}^{n} (r_i - r)^2 \leq 0.
\end{equation}
由 LaSalle 不变性原理 (LaSalle's invariance principle)，当 $t \to \infty$ 时有 $r_i \to r$。

\textbf{步骤 2：球面上的切向动力学 (Tangential dynamics on the sphere)。}
一旦对所有 $i$ 有 $r_i = r$，智能体被约束在以 $p_0$ 为中心的球面 $\mathrm{S}^{(N-1)}$ 上。
定义球面上的平均位置：
\begin{equation}
\bar{x} = \frac{1}{n} \sum_{i=1}^{n} x_i,
\end{equation}
其中 $x_i = r \sgn(e_i)$。
$\bar{x}$ 的动力学可由 \eqref{eq: error_dynamics} 推导为：
\begin{equation}
\dot{\bar{x}} = \frac{1}{n} \sum_{i=1}^{n} \dot{x}_i = \frac{1}{n} \sum_{i=1}^{n} \frac{d}{dt}(r \sgn(e_i)).
\end{equation}
由于 $r_i = r$ 为常数，$\dot{x}_i$ 完全切于球面。
由智能体间作用力可得：
\begin{equation}
\dot{\bar{x}} = \frac{k_a}{n} \sum_{i=1}^{n} \sum_{j=1}^{n} d_{ij} - \frac{k_r}{n} \sum_{i=1}^{n} \sum_{j=1, j \neq i}^{n} \frac{d_{ij}}{\|x_i - x_j\|^2}.
\label{eq: average_dynamics}
\end{equation}

\textbf{步骤 3：收敛到均匀部署 (Convergence to uniform deployment)。}
对切向动力学定义候选李雅普诺夫函数：
\begin{equation}
V_2 = \frac{1}{2} \|\bar{x}\|^2.
\end{equation}
由 \eqref{eq: average_dynamics} 及文献 [1] 中的恒等式：
\begin{equation}
\sum_{i=1}^{n} \sum_{j=1, j \neq i}^{n} \frac{d_{ij}}{\|x_i - x_j\|^2} = \frac{n(n-1)}{2r^2} \bar{x},
\end{equation}
可得：
\begin{equation}
\dot{\bar{x}} = k_a n \bar{x} - \frac{k_a}{r^2} \sum_{i=1}^{n} \langle x_i, \bar{x} \rangle x_i - \frac{k_r(n-1)}{2r^2} \bar{x}.
\end{equation}
于是：
\begin{equation}
\dot{V}_2 = \bar{x}^T \dot{\bar{x}} = k_a n \|\bar{x}\|^2 - \frac{k_a}{r^2} \sum_{i=1}^{n} \langle x_i, \bar{x} \rangle^2 - \frac{k_r(n-1)}{2r^2} \|\bar{x}\|^2.
\end{equation}
令 $\vartheta_i$ 为 $x_i$ 与 $\bar{x}$ 之间的夹角（当 $\bar{x} \neq 0$ 时），并定义 $\lambda(t) := \sum_{i=1}^{n} \cos^2 \vartheta_i$，其中 $0 < \lambda(t) \leq n$。
则：
\begin{equation}
\dot{V}_2 = k_a \left(n - \lambda(t) - \frac{(n-1)k_r}{2r^2 k_a}\right) \|\bar{x}\|^2.
\end{equation}
在增益条件 \eqref{eq:gain_condition} 下：
\begin{equation}
\frac{k_r}{k_a} \geq \frac{2n r^2}{n-1} \quad \Rightarrow \quad n - \frac{(n-1)k_r}{2r^2 k_a} \leq 0.
\end{equation}
由于 $\lambda(t) > 0$，当 $\bar{x} \neq 0$ 时有 $\dot{V}_2 < 0$，这表明当 $t \to \infty$ 时 $\bar{x} \to 0$。

\textbf{步骤 4：结论 (Conclusion)。}
由步骤 1，对所有 $i$ 有 $r_i \to r$，即 $\|p_i - p_0\| \to r$。
由步骤 3，$\bar{x} \to 0$，即 $\frac{1}{n} \sum_{i=1}^{n} x_i \to 0$。
由于 $x_i = r \sgn(p_i - p_0)$，有：
\begin{equation}
\|\bar{p} - p_0\| = \left\| \frac{1}{n} \sum_{i=1}^{n} p_i - p_0 \right\| = \left\| \frac{1}{n} \sum_{i=1}^{n} e_i \right\| = \left\| \frac{1}{n} \sum_{i=1}^{n} \frac{r_i}{r} x_i \right\| \to \left\| \frac{1}{n} \sum_{i=1}^{n} x_i \right\| = \|\bar{x}\| \to 0.
\end{equation}
因此，当 $t \to \infty$ 时有 $\|\bar{p} - p_0\| \to 0$。

证毕。
\end{proof}

\begin{theorem}[碰撞避免 (Collision Avoidance)]
\begin{equation}
\|p_i(t) - p_j(t)\| \geq \epsilon > 0, \quad \forall t \geq 0, \quad \forall i \neq j.
\end{equation}
\end{theorem}

\begin{proof}
我们采用反证法 (contradiction)，并借助障碍函数法 (method of barrier functions)。

\textbf{步骤 1：定义障碍函数 (Define barrier function)。}
考虑复合李雅普诺夫函数：
\begin{equation}
V = V_1 + V_2 + \alpha V_{\text{coll}},
\end{equation}
其中 $V_1 = \sum_{i=1}^{n} \frac{1}{2}(r_i - r)^2$ 为径向跟踪误差，$V_2 = \frac{1}{2} \|\bar{x}\|^2$ 为部署误差 (deployment error)，$\alpha > 0$ 是一个待定的充分大常数。

\textbf{步骤 2：计算时间导数 (Compute time derivative)。}
$V_{\text{coll}}$ 沿系统轨线的导数为：
\begin{equation}
\dot{V}_{\text{coll}} = \sum_{i=1}^{n} \nabla_{x_i} V_{\text{coll}}^T \dot{x}_i.
\end{equation}

由于 $\dot{x}_i$ 由控制律确定，有：
\begin{equation}
\dot{V}_{\text{coll}} = \sum_{i=1}^{n} \left( -2 \sum_{j \in \mathcal{N}_i} \frac{w_{ij} k_r}{\|x_i - x_j\|^4} (x_i - x_j) \right)^T \dot{x}_i.
\end{equation}

\textbf{步骤 3：分析主导项 (Analyze dominant term)。}
当智能体相互靠近，即 $\|x_i - x_j\| \to 0$ 时，项 $\frac{1}{\|x_i - x_j\|^4}$ 占主导。
智能体 $i$ 与 $j$ 之间的交互作用力包含：
\begin{equation}
f_{ij} = \left(k_a - \frac{k_r}{\|x_i - x_j\|^2}\right) d_{ij}.
\end{equation}
当 $\|x_i - x_j\|$ 很小时，排斥项占主导：$f_{ij} \approx -\frac{k_r}{\|x_i - x_j\|^2} d_{ij}$。

由于 $d_{ij} = x_j - \frac{\langle x_i, x_j \rangle}{r^2} x_i$ 且 $\|d_{ij}\| \leq \|x_i - x_j\|$，有：
\begin{equation}
\|f_{ij}\| \geq \frac{k_r}{\|x_i - x_j\|^2} \|d_{ij}\| \geq \frac{k_r}{\|x_i - x_j\|}.
\end{equation}

\textbf{步骤 4：界定导数 (Bound derivative)。}
排斥力将智能体推开，因此对于充分靠近的智能体：
\begin{equation}
\dot{V}_{\text{coll}} \geq \sum_{i=1}^{n} \sum_{j \in \mathcal{N}_i} \frac{c_1}{\|x_i - x_j\|^3} - c_2,
\end{equation}
其中 $c_1, c_2 > 0$ 为依赖于控制参数的常数。

\textbf{步骤 5：反证推理 (Contradiction argument)。}
假设存在一对 $(i,j)$ 和某个时刻 $t^*$，使得对任意小的 $\delta > 0$ 有 $\|x_i(t^*) - x_j(t^*)\| = \delta$。
则 $V_{\text{coll}} \geq \frac{c_1}{\delta^2}$ 可以任意大。

对充分小的 $\delta$，有：
\begin{equation}
\dot{V} \geq -k_o \sum_{i=1}^{n} (r_i - r)^2 - k_a \sum_{i=1}^{n} \|\bar{x}\|^2 + \alpha \left( \frac{c_3}{\delta^3} - c_4 \right),
\end{equation}
其中 $c_3, c_4 > 0$。

由于 $\alpha$ 可取任意大，对充分小的 $\delta$，有 $\dot{V} > 0$，这与定理 1 所示之 $V$ 沿轨线非增（$\dot{V} \leq 0$）的事实相矛盾。

因此，存在 $\epsilon > 0$，使得对所有 $t \geq 0$ 及所有 $i \neq j$，有 $\|x_i(t) - x_j(t)\| \geq \epsilon$。

\textbf{步骤 6：关联至实际位置 (Relate to actual positions)。}
由于 $x_i = r \sgn(p_i - p_0)$ 且 $x_j = r \sgn(p_j - p_0)$，有：
\begin{equation}
\|p_i - p_j\| \geq \|x_i - x_j\| \geq \epsilon,
\end{equation}
其中第一个不等式由三角不等式 (triangle inequality) 及 $\|p_i - p_0\|, \|p_j - p_0\| \to r$ 得出。

证毕。
\end{proof}



\section{模式的几何刻画 (Geometric Characterization of Patterns)}

\begin{proposition}[距离矩阵 (Distance Matrix)]
对于球形围栏控制，距离矩阵 (distance matrix) $D_n(c) := [e_{ij}] \in \mathbb{R}^{n \times n}$（其中 $e_{ij} := \lim_{t \to \infty} \|x_i - x_j\|$）满足：
\begin{equation}
\sum_{i=1}^{n-1} \sum_{j=i+1}^{n} \|x_i - x_j\|^2 = (r^2 - \|\bar{x}\|^2)n^2.
\end{equation}
当 $\bar{x} = 0$ 时，上式化为：
\begin{equation}
\sum_{i=1}^{n-1} \sum_{j=i+1}^{n} \|x_i - x_j\|^2 = r^2 n^2.
\end{equation}
\end{proposition}

\begin{proof}
对球面上的智能体，$\|x_i - x_j\|^2 = 2r^2 - 2\langle x_i, x_j \rangle$。
对所有配对求和：
\begin{equation}
\sum_{i=1}^{n} \sum_{j=1}^{n} \|x_i - x_j\|^2 = 2r^2 n^2 - 2\left\| \sum_{i=1}^{n} x_i \right\|^2 = 2r^2 n^2 - 2n^2 \|\bar{x}\|^2.
\end{equation}
注意到 $\sum_{i=1}^{n} \sum_{j=1}^{n} \|x_i - x_j\|^2 = 2 \sum_{i=1}^{n-1} \sum_{j=i+1}^{n} \|x_i - x_j\|^2$，即得结论。
\end{proof}

\begin{definition}[绝对对称模式 (Absolute Symmetric Pattern)]
若每个智能体在距离矩阵 $D_n$ 中具有相同的地位，则称该模式为\textit{绝对对称 (absolute symmetric)} 模式。
绝对对称模式的必要条件是 $\bar{x} = 0$。
\end{definition}

\section{示例说明 (Illustrative Examples)}

\begin{remark}[常见模式 (Common Patterns)]
对特定的 $n$ 值，可出现以下绝对对称模式：
\begin{itemize}
    \item $n = 3$：大圆上的等边三角形 (Equilateral triangle on a great circle)。
    \item $n = 4$：正四面体 (Regular tetrahedron)。
    \item $n = 6$：正八面体 (Regular octahedron)。
    \item $n = 8$：非立方体，而是一种更稳定的双平行正方形模式。
    \item $n = 12$：正二十面体顶点 (Regular icosahedron vertices)。
\end{itemize}
这些模式是对应智能体数量下，球面上分布最均匀的构型 (configurations)。
\end{remark}

\end{document}
